%\begin{center}
%\large \bf \runtitle
%\end{center}
%\vspace{1cm}
\chapter*{\runtitle}

\noindent This thesis studies the design and implementation of a \zkBridge for \Cardano. The goal is to authenticate an event that occurred on a source chain and enable a corresponding action on a destination chain, without relying on a centralized operator and while respecting the actual limits of \Onchain execution. To this end, the work proposes an architecture that separates an update layer, responsible for maintaining authenticated remote state, from an application layer, responsible for consuming that evidence to authorize the $\LockTx \rightarrow \MintTx$ flow.

The construction avoids verifying the full \Cardano consensus protocol directly inside a circuit. Instead, it uses \Mithril as a state-certification mechanism, a two-phase \STM verification implemented in Aiken, and auxiliary \Groth arguments for transaction membership and anti-replay \SMT updates. The prototype shows that this composition can be expressed within the \eUTxO model through canonical \UTxOs, state NFTs, \ReferenceScripts and \ProofReceipt, while remaining within the available execution-unit budgets.

The main result is a viability proof for an isomorphic \Cardano instance: not a complete production-ready heterogeneous bridge, but a verifiable architecture that precisely states its security assumptions, its \Onchain invariants and the components left for future work.

\bigskip

\noindent\textbf{Keywords:} \zkBridge, \Cardano, \Mithril, \eUTxO, zero-knowledge proofs.
